\begin{definition}[Equivariant \v Cech complex for a closed covering]
\label{definition-equivariant-cech-closed-covering}
\begin{reference}
\cite[Chapter V, \S 5.5, pp. 208--209]{Tohoku}
\end{reference}
Let $G$ act on a topological space $X$, and let
$\mathcal U=(U_i)_{i\in I}$ be a closed $G$-covering.  Thus every $U_i$ is
closed, $gU_i=U_{gi}$, and the $U_i$ cover $X$.  Put
$U_{i_0\ldots i_p}=U_{i_0}\cap\cdots\cap U_{i_p}$.  If $\mathcal P$ is an
abelian $G$-equivariant coefficient system on these finite intersections,
write
$$
C^p_{\mathrm{cl}}(\mathcal U,\mathcal P)=
\prod_{(i_0,\ldots,i_p)\in I^{p+1}}
\mathcal P(U_{i_0\ldots i_p}).
$$
Here $\mathcal P(U_{\mathbf i})$ means
$\Gamma(U_{\mathbf i},\mathcal P|_{U_{\mathbf i}})$ when $\mathcal P$ is a
sheaf; the coefficient-system transport maps carry this group to the group
on $gU_{\mathbf i}$.  The alternating restriction maps make this a complex
of $G$-modules.  We define
$$
\check H^n_{\mathrm{cl}}(\mathcal U;G,\mathcal P)=
H^n\left(R\Gamma_G\bigl(
C^\bullet_{\mathrm{cl}}(\mathcal U,\mathcal P)\bigr)\right).
$$
For a $G$-sheaf $\mathcal A$, the coefficient system used below is
$$
\underline H^q_{\mathrm{cl}}(\mathcal A)(U_{i_0\ldots i_p})
 =H^q(U_{i_0\ldots i_p},\mathcal A|_{U_{i_0\ldots i_p}}),
$$
with transport induced by the action.  This is the same derived-invariants
construction as Definition \ref{definition-equivariant-cech-cohomology};
the subscript only records that the covering members need not be open.
\end{definition}

\begin{lemma}[Equivariant \v Cech spectral sequence for closed coverings]
\label{lemma-equivariant-cech-spectral-sequence-closed-covering}
\begin{reference}
\cite[Chapter V, \S 5.5, Theorem 5.5.3 and Corollary 2,
pp. 208--209]{Tohoku}
\end{reference}
Let $G$ act on a topological space $X$, let
$\mathcal U=(U_i)_{i\in I}$ be a closed $G$-covering, and let $\mathcal A$ be
an abelian $G$-sheaf on $X$.  Assume either that $I$ is finite, or that $X$
is paracompact and $\mathcal U$ is locally finite.  There is a first-quadrant
spectral sequence, additive and functorial in $\mathcal A$ from the
$E_2$-page onward,
$$
E_2^{p,q}=\check H^p_{\mathrm{cl}}
 (\mathcal U;G,\underline H^q_{\mathrm{cl}}(\mathcal A))
 \Longrightarrow H^{p+q}(X;G,\mathcal A),
$$
with a finite filtration in each total degree.  Its bottom edge maps are
canonical functorial homomorphisms
$$
\check H^n_{\mathrm{cl}}(\mathcal U;G,\mathcal A)
 \longrightarrow H^n(X;G,\mathcal A).
$$
They are isomorphisms for $n=0$ and injective for $n=1$, and the associated
five-term exact sequence begins
$$
\begin{aligned}
0\longrightarrow{}&\check H^1_{\mathrm{cl}}(\mathcal U;G,\mathcal A)
\longrightarrow H^1(X;G,\mathcal A)
\longrightarrow \check H^0_{\mathrm{cl}}
 (\mathcal U;G,\underline H^1_{\mathrm{cl}}(\mathcal A))\\
\longrightarrow{}&\check H^2_{\mathrm{cl}}(\mathcal U;G,\mathcal A)
\longrightarrow H^2(X;G,\mathcal A).
\end{aligned}
$$
\end{lemma}

\begin{proof}
For a tuple $\mathbf i=(i_0,\ldots,i_r)$ let
$j_{\mathbf i}:U_{\mathbf i}\hookrightarrow X$ be the closed inclusion and
put
$$
\mathcal K^r(\mathcal U,\mathcal A)=
\prod_{\mathbf i\in I^{r+1}}
(j_{\mathbf i})_*(\mathcal A|_{U_{\mathbf i}}),
$$
with the alternating restriction differential.  The augmentation
$$
\mathcal A[0]\longrightarrow
\mathcal K^\bullet(\mathcal U,\mathcal A)
$$
is equivariant, where $G$ acts by transport and permutation of the factors.
It is a quasi-isomorphism.  Indeed, the stalk argument in the proof of Lemma
\ref{lemma-cech-spectral-sequence-closed-covering} applies unchanged: in the
finite case, and locally in the locally finite case, the nonzero factors are
the augmented cochain complex of a nonempty simplex.

Let $F$ be global sections from abelian $G$-sheaves to $G$-modules and let
$G_0$ be the invariant-module functor.  If $\mathcal I$ is an injective
abelian $G$-sheaf, then $F(\mathcal I)$ is right acyclic for $G_0$ by
Injectives, Lemma
\ref{injectives-lemma-equivariant-injective-sections-invariants}.  Apply
Derived Categories, Lemma
\ref{derived-lemma-hyperderived-composite-functor-spectral-sequence-resolution}
to the displayed equivariant resolution.  Its $E_2$-page is
$$
H^p\left(RG_0\bigl(
(R^qF)(\mathcal K^\bullet(\mathcal U,\mathcal A))\bigr)\right)
$$
and it converges to $R^{p+q}(G_0\circ F)(\mathcal A)$.

It remains to identify the complex inside $RG_0$.  The cohomology
calculation in the proof of Lemma
\ref{lemma-cech-spectral-sequence-closed-covering} is functorial and
equivariant.  Injectives, Lemma
\ref{injectives-lemma-equivariant-injective-underlying} identifies the
underlying groups of $R^qF$ with ordinary sheaf cohomology, carrying the
induced $G$-action.  In the finite case the closed-cover calculation uses
exact direct image and finite products for closed inclusions; in the locally
finite paracompact case it uses the soft resolution and locally finite
products.  It gives
$$
(R^qF)(\mathcal K^\bullet(\mathcal U,\mathcal A))
\cong
C^\bullet_{\mathrm{cl}}(
\mathcal U,\underline H^q_{\mathrm{cl}}(\mathcal A))
$$
as complexes of $G$-modules.  Definition
\ref{definition-equivariant-cech-closed-covering} now identifies the
$E_2$-page with the one in the statement, while Definition
\ref{definition-orbit-space-equivariant-cohomology} identifies the abutment
with $H^{p+q}(X;G,\mathcal A)$.  The general edge and convergence statements
give the finite filtrations, the bottom edge maps, and the five-term exact
sequence asserted above.
\end{proof}

\begin{lemma}
\label{lemma-equivariant-cech-spectral-sequence-closed-acyclic-cover}
\begin{reference}
\cite[Chapter V, \S 5.5, Corollaries 1 and 2 to Theorem 5.5.3,
pp. 208--209]{Tohoku}
\end{reference}
In the situation of Lemma
\ref{lemma-equivariant-cech-spectral-sequence-closed-covering}, assume in
addition that
$$
H^q(U_{i_0\ldots i_p},\mathcal A|_{U_{i_0\ldots i_p}})=0
\quad\text{for every }q>0\text{ and every }(i_0,\ldots,i_p).
$$
Then all bottom edge maps
$$
\check H^n_{\mathrm{cl}}(\mathcal U;G,\mathcal A)
\longrightarrow H^n(X;G,\mathcal A)
$$
are isomorphisms.
\end{lemma}

\begin{proof}
The $E_2$-page of Lemma
\ref{lemma-equivariant-cech-spectral-sequence-closed-covering} is
concentrated in the row $q=0$.  The spectral sequence therefore degenerates,
and its bottom edge maps are isomorphisms in every degree.
\end{proof}
